 >>?% ======================================================================
% CAPITULO 7 --- Conexao, Transporte Paralelo e Derivadas de Nelson
% Baseado em: Farinelli, S. (2021). arXiv:0910.1671v10, Secoes 3.4--3.6
% ======================================================================
\chapter{Conex  o, Transporte Paralelo e Derivadas de Nelson}

Nos dois cap -tulos anteriores, constru -mos as pe SSas do tabuleiro: os
\textit{gauges} (Cap -tulo~5) como objetos que codificam a informa SS  o
financeira de um instrumento, e o \textit{fibrado principal do mercado}
(Cap -tulo~6) como o espa SSo no qual estes objetos vivem.  Falta a pe SSa
central: a \textbf{conex  o} --- a regra que nos permite ``transportar''
informa SS  o financeira ao longo de curvas no espa SSo de carteiras e no
tempo.

Este cap -tulo  c o cora SS  o operacional da GAT.  Definiremos a conex  o de
Ilinski--Farinelli, provaremos que o transporte paralelo corresponde a
opera SS ues financeiras fundamentais (c  mbio e desconto), e
introduziremos a derivada de Nelson --- a ferramenta anal -tica que
permite formular o autofinanciamento na linguagem geom ctrica.

% ----------------------------------------------------------------------
\section{O Que  %% uma Conex  o?}
% ----------------------------------------------------------------------

\subsection{Intui SS  o: Conectando Fibras Vizinhas}

Considere o fibrado principal $\mathcal{B}$ sobre a base $M = X \times
[0,+\infty[$ que constru -mos no Cap -tulo~6.  Cada ponto $p = (x,t) \in
M$ tem uma fibra $\mathcal{B}_p$ isomorfa ao grupo $G^*$ (as
transforma SS ues de gauge invert -veis).  Os elementos da fibra s  o
diferentes representa SS ues do mesmo estado financeiro --- diferentes
escolhas de num crario, diferentes maneiras de empacotar fluxos de
caixa.

O problema fundamental  c: \textbf{como comparar elementos em fibras
vizinhas?}  Se eu tenho um gauge em $(x,t)$ e quero saber qual gauge
em $(x+\delta x, t+\delta t)$ lhe corresponde, preciso de uma regra.
Esta regra  c a \textbf{conex  o}.

\begin{defi}[Conex  o de Ehresmann --- Decomposi SS  o do Espa SSo Tangente]\label{def:ehresmann}
Seja $\pi: \mathcal{B} \to M$ um fibrado principal com grupo estrutural
$G$.  Uma \textbf{conex  o} (de Ehresmann) em $\mathcal{B}$  c uma
distribui SS  o suave de subespa SSos horizontais $H\mathcal{B} \subset
T\mathcal{B}$ tal que, para cada $b \in \mathcal{B}$:
\begin{enumerate}[(i)]
  \item $T_b\mathcal{B} = V_b\mathcal{B} \oplus H_b\mathcal{B}$
        (decomposi SS  o em soma direta);
  \item $H_{b\cdot g}\mathcal{B} = (R_g)_* H_b\mathcal{B}$
        (equivari  ncia sob a a SS  o    direita de $G$);
  \item $H_b\mathcal{B}$ depende suavemente de $b$.
\end{enumerate}
Aqui, $V_b\mathcal{B} := \ker(d\pi_b)$  c o \textbf{espa SSo vertical}
(tangente    fibra) e $H_b\mathcal{B}$  c o \textbf{espa SSo horizontal}
(complementar, isomorfo a $T_{\pi(b)}M$ via $d\pi_b$).
\end{defi}

\begin{obs}[Interpreta SS  o Geom ctrica]
O espa SSo vertical $V_b\mathcal{B}$ cont cm as dire SS ues que permanecem
na mesma fibra --- transforma SS ues de gauge que n  o alteram o ponto-base
$(x,t)$.  O espa SSo horizontal $H_b\mathcal{B}$ cont cm as dire SS ues que
nos levam a pontos-base vizinhos --- mudan SSas de carteira ou avan SSo no
tempo.  A conex  o decide, para cada dire SS  o na base, qual dire SS  o na
fibra  c considerada ``paralela''.
\end{obs}

\subsection{Proje SS ues e a Forma de Conex  o}

A decomposi SS  o $T\mathcal{B} = V\mathcal{B} \oplus H\mathcal{B}$
define dois projetores:
\begin{equation}\label{eq:projectors}
  \Pi^v: T\mathcal{B} \to V\mathcal{B}, \qquad
  \Pi^h: T\mathcal{B} \to H\mathcal{B},
\end{equation}
com $\Pi^v + \Pi^h = \operatorname{id}$ e $(\Pi^v)^2 = \Pi^v$,
$(\Pi^h)^2 = \Pi^h$.

No fibrado principal, a estrutura de grupo permite codificar a conex  o
de forma mais eficiente atrav cs da \textbf{forma de conex  o}: uma
1-forma $\chi$ na base $M$ com valores na  ilgebra de Lie $\mathfrak{g}$
do grupo $G$.  Dado um vetor tangente $v \in T_pM$, a forma $\chi_p(v)$
prescreve o elemento infinitesimal de $\mathfrak{g}$ pelo qual devemos
multiplicar o elemento da fibra para permanecermos na horizontal.

% ----------------------------------------------------------------------
\section{A Conex  o de Mercado (Ilinski--Farinelli)}
% ----------------------------------------------------------------------

\subsection{Motiva SS  o Financeira}

Queremos que o transporte paralelo tenha uma interpreta SS  o financeira
transparente.  H i duas dire SS ues fundamentais na base:

\begin{itemize}
  \item \textbf{Dire SS ues nominais} ($\delta x$, com $\delta t = 0$):
        mudar de carteira no mesmo instante.  Duas carteiras $x$ e
        $x + \delta x$ que diferem apenas na composi SS  o de ativos
        est  o relacionadas por uma \textbf{taxa de c  mbio}: quantas
        unidades da carteira $x$ equivalem a uma unidade da carteira
        $x + \delta x$?  A resposta  c $D_t^{x+\delta x} / D_t^x$.
  \item \textbf{Dire SS  o temporal} ($\delta t$, com $\delta x = 0$):
        avan SSar no tempo com a mesma carteira.  O valor de uma
        carteira hoje e amanh   est  o relacionados pelo \textbf{fator
        de desconto}: $1$ unidade hoje equivale a $e^{\int_t^{t+\delta
        t} r_u^x du}$ unidades em $t+\delta t$.
\end{itemize}

A conex  o deve capturar ambas as rela SS ues em uma  onica express  o.

\subsection{Defini SS  o da Conex  o}

\begin{defi}[Conex  o de Ilinski--Farinelli]\label{def:market_connection}
A \textbf{conex  o de mercado} no fibrado $\mathcal{B}$  c definida pela
prescri SS  o do transporte paralelo ao longo de curvas diferenci iveis.
Para um deslocamento infinitesimal $(\delta x, \delta t)$ a partir de
$(x,t)$, o elemento do grupo $G^*$ que implementa o levantamento
horizontal  c:

\begin{equation}\label{eq:connection_definition}
  \Gamma\big((x,t),(\delta x, \delta t)\big) \cdot g
  := g\left(\frac{D_t^{x+\delta x}}{D_t^x} - r_t^x\,\delta t\right).
\end{equation}

Equivalentemente, a \textbf{forma de conex  o} $\chi$ --- uma 1-forma
na base $M$ com valores na  ilgebra de Lie $\mathfrak{g} \cong
\mathbb{R}^{[0,+\infty[}$ ---  c dada por:

\begin{equation}\label{eq:connection_form}
  \chi(x,t)(\delta x, \delta t)
  = \left(\frac{D_t^{\delta x}}{D_t^x} - r_t^x\,\delta t\right).
\end{equation}
\end{defi}

\begin{obs}[Linearidade em $\delta x$]
A nota SS  o $D_t^{\delta x}$ merece esclarecimento.  Para uma varia SS  o
infinitesimal $\delta x = (\delta x^1, \dots, \delta x^N)$, o deflator
da carteira $\delta x$  c a combina SS  o linear:
\begin{equation}\label{eq:D_delta_x}
  D_t^{\delta x} := \sum_{j=1}^N D_t^j \cdot \delta x^j,
\end{equation}
onde $D_t^j$  c o deflator do $j$- csimo ativo (o valor presente de uma
perpetuidade unit iria associada   quele ativo).  Para a carteira
$x + \delta x$, temos $D_t^{x+\delta x} = D_t^x + D_t^{\delta x}$, de
modo que:
\[
\frac{D_t^{x+\delta x} - D_t^x}{D_t^x}
= \frac{D_t^{\delta x}}{D_t^x}.
\]
Esta  c precisamente a taxa de varia SS  o relativa do deflator quando
mudamos de $x$ para $x + \delta x$ --- a \textbf{taxa de c  mbio
infinitesimal} entre as duas carteiras.
\end{obs}

\subsection{Express  o em Coordenadas}

Em coordenadas locais $(x^1, \dots, x^N, t)$ na base $M$, a forma de
conex  o $\chi$ escreve-se como:

\begin{equation}\label{eq:connection_coordinates}
  \chi(x,t) = \frac{1}{D_t^x}\sum_{j=1}^N D_t^j\,dx^j \;-\; r_t^x\,dt.
\end{equation}

Verifica SS  o: aplicando a um vetor $(\delta x, \delta t) = (\delta x^1,
\dots, \delta x^N, \delta t)$, obtemos:
\begin{align*}
  \chi(x,t)(\delta x, \delta t)
  &= \frac{1}{D_t^x}\sum_{j=1}^N D_t^j\,\delta x^j \;-\; r_t^x\,\delta t \\
  &= \frac{D_t^{\delta x}}{D_t^x} - r_t^x\,\delta t,
\end{align*}
que coincide com~\eqref{eq:connection_form}.

\begin{obs}[A Taxa Curta da Carteira $x$]
A quantidade $r_t^x$ que aparece na conex  o  c a \textbf{taxa curta}
(ou \textit{short rate}) associada    carteira $x$, definida como:
\begin{equation}\label{eq:short_rate_x}
  r_t^x := -\left.\frac{\partial}{\partial s}\log P_{t,s}^x\right|_{s=t},
\end{equation}
onde $P_{t,s}^x$  c a estrutura a termo da carteira $x$.  Em termos do
deflator:
\begin{equation}\label{eq:short_rate_deflator}
  r_t^x = -\mathcal{D}\log D_t^x,
\end{equation}
com $\mathcal{D}$ denotando a derivada m cdia de Nelson (derivada de
Stratonovich), que ser i formalizada na Se SS  o~7.6.

Note que $r_t^x$ \textbf{n  o  c linear} em $x$ --- a taxa curta de uma
combina SS  o de carteiras n  o  c a combina SS  o das taxas curtas.  Esta
n  o-linearidade  c precisamente a fonte de \textbf{curvatura} (e,
portanto, de \textbf{arbitragem}) na GAT.
\end{obs}

% ----------------------------------------------------------------------
\section{A Forma de Conex  o $\chi$ como 1-Forma com Valores em $\mathfrak{g}$}
% ----------------------------------------------------------------------

\subsection{  lgebra de Lie do Grupo de Gauge}

O grupo estrutural  c $G^*$, o grupo Abeliano das transforma SS ues de
gauge invert -veis (intensidades de cashflow com inversa sob
convolu SS  o).  Por ser Abeliano, sua  ilgebra de Lie $\mathfrak{g}$  c
isomorfa ao espa SSo vetorial $\mathbb{R}^{[0,+\infty[}$ (o espa SSo de
todas as intensidades de cashflow) com o colchete de Lie trivial
$[\cdot, \cdot] \equiv 0$.

\begin{prop}[Propriedades da Forma de Conex  o]\label{prop:chi_properties}
A 1-forma $\chi$ com valores em $\mathfrak{g}$ satisfaz:
\begin{enumerate}[(i)]
  \item $\chi$  c \textbf{horizontal}: para qualquer vetor vertical
        $v \in V_b\mathcal{B}$, $\chi(v) = 0$.
  \item $\chi$  c \textbf{equivariante}: para $g \in G^*$,
        $(R_g)^*\chi = \operatorname{Ad}_{g^{-1}} \circ \chi$.
        Como $G^*$  c Abeliano, $\operatorname{Ad}_g = \operatorname{id}$,
        logo $(R_g)^*\chi = \chi$ (invari  ncia    direita).
  \item $\chi$  c \textbf{suave} nos argumentos $(x,t)$.
\end{enumerate}
\end{prop}

\subsection{Express  o Completa com a Fibra}

Quando inclu -mos a coordenada da fibra $g \in G^*$, a forma de conex  o
$\omega$ no fibrado total $\mathcal{B}$ escreve-se:

\begin{equation}\label{eq:connection_form_full}
  \omega(x,t,g)(\delta x, \delta t, \delta g)
  = \delta g \cdot g^{-1} + \operatorname{Ad}_g \chi(x,t)(\delta x, \delta t).
\end{equation}

Como $G^*$  c Abeliano, $\operatorname{Ad}_g = \operatorname{id}$, e
esta express  o simplifica-se para:

\begin{equation}\label{eq:connection_form_simple}
  \omega(x,t,g)(\delta x, \delta t, \delta g)
  = \delta g \cdot g^{-1} + \chi(x,t)(\delta x, \delta t).
\end{equation}

O primeiro termo $\delta g \cdot g^{-1}$  c a \textbf{forma de
Maurer--Cartan} de $G^*$, que mede varia SS ues puramente verticais
(dentro da fibra).  O segundo termo $\chi(x,t)(\delta x, \delta t)$  c a
contribui SS  o horizontal, que depende apenas do ponto-base.

% ----------------------------------------------------------------------
\section{Transporte Paralelo ao Longo de Curvas}
% ----------------------------------------------------------------------

\subsection{Levantamento Horizontal}

Dada uma curva suave $\gamma: [0,1] \to M$ na base, com
$\gamma(s) = (x(s), t(s))$, e um ponto inicial $g_0 \in
\mathcal{B}_{\gamma(0)}$ na fibra sobre $\gamma(0)$, o
\textbf{levantamento horizontal} $\tilde{\gamma}: [0,1] \to
\mathcal{B}$  c a  onica curva no fibrado tal que:

\begin{enumerate}[(i)]
  \item $\pi \circ \tilde{\gamma} = \gamma$ (projeta-se sobre $\gamma$);
  \item $\tilde{\gamma}(0) = g_0$ (condi SS  o inicial);
  \item $\tilde{\gamma}'(s) \in H_{\tilde{\gamma}(s)}\mathcal{B}$ para
        todo $s$ (a curva  c horizontal).
\end{enumerate}

Em termos da forma de conex  o, a condi SS  o de horizontalidade  c:

\begin{equation}\label{eq:horizontality_condition}
  \omega\big(\tilde{\gamma}(s)\big)\big(\tilde{\gamma}'(s)\big) = 0,
  \quad \forall s \in [0,1].
\end{equation}

Usando~\eqref{eq:connection_form_simple}, obtemos a equa SS  o diferencial
para o levantamento:

\begin{equation}\label{eq:horizontal_lift_ode}
  \frac{dg}{ds} \cdot g^{-1} + \chi\big(x(s), t(s)\big)\big(x'(s), t'(s)\big) = 0.
\end{equation}

Como $G^*$  c Abeliano e escrevemos a a SS  o como multiplica SS  o, esta
equa SS  o tem a solu SS  o expl -cita:

\begin{equation}\label{eq:parallel_transport_solution}
  g(s) = g_0 \cdot \exp\left(-\int_0^s
         \chi\big(x(u), t(u)\big)\big(x'(u), t'(u)\big)\,du\right).
\end{equation}

\subsection{Defini SS  o do Transporte Paralelo}

\begin{defi}[Transporte Paralelo]\label{def:parallel_transport}
Seja $\gamma: [0,1] \to M$ uma curva suave na base.  O
\textbf{transporte paralelo} ao longo de $\gamma$  c o isomorfismo
entre fibras:
\begin{equation}\label{eq:parallel_transport_map}
  \Pi_\gamma: \mathcal{B}_{\gamma(0)} \to \mathcal{B}_{\gamma(1)},
  \quad
  \Pi_\gamma(g_0) := g_0 \cdot \exp\left(-\int_\gamma \chi\right),
\end{equation}
onde $\int_\gamma \chi$ denota a integral de linha da 1-forma $\chi$
ao longo de $\gamma$:
\begin{equation}\label{eq:line_integral_chi}
  \int_\gamma \chi := \int_0^1
  \chi\big(\gamma(s)\big)\big(\gamma'(s)\big)\,ds.
\end{equation}
\end{defi}

\begin{obs}[Estrutura de Grupo do Transporte Paralelo]
O transporte paralelo  c um \textbf{isomorfismo de $G^*$-espa SSos}:
$\Pi_\gamma(g \cdot h) = \Pi_\gamma(g) \cdot h$ para todo $h \in
G^*$ (equivari  ncia).  Al cm disso:
\begin{itemize}
  \item $\Pi_{\gamma^{-1}} = \Pi_\gamma^{-1}$ (reversibilidade);
  \item $\Pi_{\gamma_2 \ast \gamma_1} = \Pi_{\gamma_2} \circ
        \Pi_{\gamma_1}$ (composi SS  o de caminhos).
\end{itemize}
\end{obs}

% ----------------------------------------------------------------------
\section{Interpreta SS  o Financeira do Transporte Paralelo}
% ----------------------------------------------------------------------

A beleza da constru SS  o de Ilinski--Farinelli reside em que o transporte
paralelo --- um conceito puramente geom ctrico --- admite uma
interpreta SS  o financeira exata e intuitiva.

\begin{teo}[Interpreta SS  o Financeira do Transporte Paralelo ---
           Teorema~29 do Artigo]\label{teo:financial_interpretation}
Seja $\Pi_\gamma$ o transporte paralelo definido pela
conex  o~\eqref{eq:connection_definition}.  Ent  o:

\begin{enumerate}[(a)]
  \item \textbf{Ao longo de $x$-linhas ($t$ fixo):}
        Seja $\gamma(s) = (x_0 + s \cdot \Delta x, t_0)$ uma curva que
        conecta duas carteiras no mesmo instante.  O transporte
        paralelo ao longo de $\gamma$ corresponde   
        \textbf{multiplica SS  o pela taxa de c  mbio} entre as carteiras:
        \begin{equation}\label{eq:parallel_transport_x}
          \Pi_\gamma(g_0) = g_0 \cdot \frac{D_{t_0}^{x_0}}{D_{t_0}^{x_0 + \Delta x}}.
        \end{equation}
        Em palavras: transportar paralelamente um gauge da carteira
        $x_0$ para a carteira $x_0 + \Delta x$ equivale a multiplic i-lo
        pela raz  o dos deflatores --- exatamente o que uma casa de
        c  mbio faz ao converter moedas.

  \item \textbf{Ao longo de $t$-linhas ($x$ fixo):}
        Seja $\gamma(s) = (x_0, t_0 + s)$ uma curva que avan SSa no
        tempo mantendo a mesma carteira.  O transporte paralelo ao
        longo de $\gamma$ corresponde    \textbf{divis  o pelo fator de
        desconto}:
        \begin{equation}\label{eq:parallel_transport_t}
          \Pi_\gamma(g_0) = g_0 \cdot \exp\left(\int_{t_0}^{t_0 + \Delta t} r_u^{x_0}\,du\right)^{-1}.
        \end{equation}
        Equivalentemente, $\Pi_\gamma(g_0)$  c o valor descontado do
        gauge $g_0$ --- trazido a valor presente pela taxa curta da
        carteira $x_0$.
\end{enumerate}
\end{teo}

\begin{proof}
Provemos cada caso separadamente.

\textbf{Caso (a) --- $x$-linhas.}  A curva  c $\gamma(s) = (x_0 + s
\Delta x, t_0)$, com $s \in [0,1]$ e $\Delta x$ fixo.  Temos:
\begin{align*}
  \gamma'(s) &= (\Delta x, 0), \\
  \chi(\gamma(s))(\gamma'(s))
  &= \frac{D_{t_0}^{\Delta x}}{D_{t_0}^{x_0 + s\Delta x}} - r_{t_0}^{x_0 + s\Delta x} \cdot 0
   = \frac{D_{t_0}^{\Delta x}}{D_{t_0}^{x_0 + s\Delta x}}.
\end{align*}
A integral de linha  c:
\[
\int_0^1 \frac{D_{t_0}^{\Delta x}}{D_{t_0}^{x_0 + s\Delta x}}\,ds
= \int_0^1 \frac{d}{ds}\log D_{t_0}^{x_0 + s\Delta x}\,ds
= \log D_{t_0}^{x_0 + \Delta x} - \log D_{t_0}^{x_0}
= \log\frac{D_{t_0}^{x_0 + \Delta x}}{D_{t_0}^{x_0}}.
\]
Pela defini SS  o de transporte paralelo~\eqref{eq:parallel_transport_map}:
\[
\Pi_\gamma(g_0) = g_0 \cdot \exp\!\left(-\log\frac{D_{t_0}^{x_0 + \Delta x}}{D_{t_0}^{x_0}}\right)
= g_0 \cdot \frac{D_{t_0}^{x_0}}{D_{t_0}^{x_0 + \Delta x}},
\]
confirmando~\eqref{eq:parallel_transport_x}.

\textbf{Caso (b) --- $t$-linhas.}  A curva  c $\gamma(s) = (x_0, t_0 +
s)$, com $s \in [0, \Delta t]$.  Temos:
\begin{align*}
  \gamma'(s) &= (0, 1), \\
  \chi(\gamma(s))(\gamma'(s))
  &= \frac{0}{D_{t_0+s}^{x_0}} - r_{t_0+s}^{x_0} \cdot 1
   = -\,r_{t_0+s}^{x_0}.
\end{align*}
A integral de linha  c:
\[
\int_0^{\Delta t} \bigl(-r_{t_0+s}^{x_0}\bigr)\,ds
= -\int_{t_0}^{t_0+\Delta t} r_u^{x_0}\,du.
\]
Portanto:
\[
\Pi_\gamma(g_0) = g_0 \cdot \exp\!\left(\int_{t_0}^{t_0+\Delta t} r_u^{x_0}\,du\right)
= g_0 \cdot \exp\!\left(\int_{t_0}^{t_0+\Delta t} r_u^{x_0}\,du\right)^{-1},
\]
confirmando~\eqref{eq:parallel_transport_t}.  O sinal negativo na
integral reflete o fato de que $1$ unidade hoje vale \textbf{mais} do
que $1$ unidade no futuro --- o fator de desconto  c o \textbf{inverso}
do fator de capitaliza SS  o.
\end{proof}

\begin{obs}[Unifica SS  o Conceitual]
O Teorema~\ref{teo:financial_interpretation} revela a profunda unidade
por tr is de duas opera SS ues financeiras aparentemente distintas:
\begin{itemize}
  \item \textbf{C  mbio} (converter entre ativos no mesmo instante) e
  \item \textbf{Desconto} (trazer valores futuros a presente)
\end{itemize}
s  o manifesta SS ues do \textbf{mesmo} fen 'meno geom ctrico --- o
transporte paralelo ao longo de diferentes dire SS ues no espa SSo de
carteiras-tempo.  A conex  o $\chi$ unifica ambas em um  onico objeto
matem itico.
\end{obs}

% ----------------------------------------------------------------------
\section{O Modelo de Mercado $D$-Diferenci ivel de Nelson}
% ----------------------------------------------------------------------

\subsection{A Derivada M cdia de Nelson}

Para formular a din  mica estoc istica do mercado na linguagem da
geometria diferencial, precisamos de uma no SS  o de derivada que seja
compat -vel com a integral de Stratonovich.  Esta  c a \textbf{derivada
m cdia de Nelson} (ou derivada estoc istica sim ctrica)\footnote{Nelson, E.
\textit{Dynamical Theories of Brownian Motion}. Princeton University
Press, 1967. ISBN: 978-0691079509. \textbf{Resenha:} Obra seminal que
reformulou o movimento Browniano como um processo estoc istico com
derivadas m cdia forward e backward, antecipando em d ccadas o c ilculo de
Malliavin e a mec  nica estoc istica.  A derivada de Nelson  c o an ilogo
estoc istico da derivada usual e satisfaz a regra de Leibniz.}.

\begin{defi}[Derivada M cdia de Nelson]\label{def:nelson_derivative}
Seja $X = (X_t)_{t \geq 0}$ um processo estoc istico adaptado.  A
\textbf{derivada m cdia de Nelson} (ou \textit{Nelson derivative}) de
$X$ no instante $t$  c definida como o limite em probabilidade:
\begin{equation}\label{eq:nelson_def}
  \mathcal{D}X_t := \lim_{\varepsilon \to 0^+}\,
  \mathbb{E}\!\left[
    \frac{X_{t+\varepsilon} - X_{t-\varepsilon}}{2\varepsilon}
    \;\middle|\;
    \mathcal{A}_t
  \right],
\end{equation}
sempre que o limite existir.  Dizemos que $X$  c \textbf{$D$-diferenci ivel}
no sentido de Nelson se $\mathcal{D}X_t$ existe para todo $t \geq 0$.
\end{defi}

\begin{obs}[Simetria Temporal]
A derivada de Nelson  c \textbf{sim ctrica no tempo}: ela utiliza
incrementos tanto para o futuro ($X_{t+\varepsilon}$) quanto para o
passado ($X_{t-\varepsilon}$).  Esta simetria  c crucial: enquanto a
derivada de It ' ($\lim_{\varepsilon \to 0} \mathbb{E}[(X_{t+\varepsilon}
- X_t)/\varepsilon \mid \mathcal{A}_t]$) olha apenas para o futuro, a
derivada de Nelson olha para ambas as dire SS ues.  O resultado  c que
\textbf{a derivada de Nelson satisfaz a regra da cadeia usual}:
\begin{equation}\label{eq:nelson_chain}
  \mathcal{D}\big(f(X_t)\big) = f'(X_t)\,\mathcal{D}X_t,
\end{equation}
para toda $f \in C^1$.  Esta propriedade  c exatamente o que torna a
integral de Stratonovich compat -vel com a geometria diferencial.
\end{obs}

\begin{obs}[Rela SS  o com Stratonovich]
Para processos de It ' da forma $dX_t = \mu_t\,dt + \sigma_t\,dW_t$,
a derivada de Nelson coincide com o \textbf{drift de Stratonovich}:
\begin{equation}\label{eq:nelson_stratonovich}
  \mathcal{D}X_t = \mu_t,
\end{equation}
enquanto o drift de It '  c $\mathbb{E}[dX_t/dt \mid \mathcal{A}_t] =
\mu_t\,dt$.  A diferen SSa entre as duas no SS ues  c precisamente o termo de
corre SS  o de It ': $\frac{1}{2}d\langle X, X\rangle_t$, que a derivada de
Nelson absorve naturalmente pela simetriza SS  o temporal.
\end{obs}

\subsection{Covaria SS  o de Nelson}

\begin{defi}[Covaria SS  o de Nelson]\label{def:nelson_covariation}
Para dois processos $D$-diferenci iveis $X$ e $Y$, define-se a
\textbf{derivada de Nelson da covaria SS  o quadr itica}:
\begin{equation}\label{eq:nelson_covar}
  \mathcal{D}\langle X, Y\rangle_t
  := \lim_{\varepsilon \to 0^+}\,
  \mathbb{E}\!\left[
    \frac{(X_{t+\varepsilon} - X_{t-\varepsilon})
          (Y_{t+\varepsilon} - Y_{t-\varepsilon})}
         {2\varepsilon}
    \;\middle|\;
    \mathcal{A}_t
  \right].
\end{equation}
Esta quantidade captura a ``volatilidade instant  nea conjunta'' dos
processos $X$ e $Y$.
\end{defi}

\begin{prop}[Regra de Leibniz para a Derivada de Nelson]
\label{prop:nelson_leibniz}
Se $X$ e $Y$ s  o $D$-diferenci iveis, ent  o o produto $XY$ tamb cm o  c e:
\begin{equation}\label{eq:nelson_leibniz}
  \mathcal{D}(X_t Y_t) = (\mathcal{D}X_t) Y_t + X_t (\mathcal{D}Y_t)
                         + \frac{1}{2}\,\mathcal{D}\langle X, Y\rangle_t.
\end{equation}
O termo de covaria SS  o $\frac{1}{2}\mathcal{D}\langle X, Y\rangle_t$  c
o \textbf{termo de It '} --- a corre SS  o de segunda ordem que distingue o
c ilculo estoc istico do c ilculo determin -stico usual.
\end{prop}

\begin{proof}[Esbo SSo]
Pela defini SS  o~\eqref{eq:nelson_def}, considere $X_{t\pm\varepsilon} =
X_t \pm \varepsilon\,\mathcal{D}X_t + o(\varepsilon)$ e analogamente
para $Y$.  Expandindo o produto $(X_{t+\varepsilon} - X_{t-\varepsilon})
(Y_{t+\varepsilon} - Y_{t-\varepsilon})$ e tomando a esperan SSa
condicional, o termo cruzado $\mathcal{D}X_t \cdot \mathcal{D}Y_t \cdot
\varepsilon^2$ contribui com a covaria SS  o no limite $\varepsilon \to 0$.
\end{proof}

\subsection{O Modelo de Mercado $D$-Diferenci ivel}

\begin{defi}[Modelo de Mercado $D$-Diferenci ivel ---
           Defini SS  o~32 do Artigo]\label{def:D_diff_market}
Um mercado  c dito \textbf{$D$-diferenci ivel} (no sentido de Nelson) se:
\begin{enumerate}[(i)]
  \item Todos os deflatores $D_t^j$ ($j = 1, \dots, N$) e o pricing
        kernel $\beta_t$ (state price deflator) s  o processos
        $D$-diferenci iveis;
  \item A taxa curta $r_t^0 = -\mathcal{D}\log S_t^0$  c
        $D$-diferenci ivel;
  \item As estruturas a termo $P_{t,s}^j$ s  o suaves em $s$ e
        $D$-diferenci iveis em $t$.
\end{enumerate}
\end{defi}

A condi SS  o de $D$-diferenciabilidade  c bastante geral: inclui todos os
processos de It ' (difus ues), processos de salto com intensidade
diferenci ivel, e combina SS ues de ambos.  Virtualmente todos os modelos
de mercado usados na pr itica (Black--Scholes, Heston, Bates, SABR,
etc.) s  o $D$-diferenci iveis.

\subsection{Estrat cgias Fracamente $D$-Diferenci iveis}

Para aplicar o c ilculo de Nelson a estrat cgias de negocia SS  o, precisamos
de uma no SS  o adequada de diferenciabilidade para processos previs -veis.

\begin{defi}[Estrat cgia Fracamente $D$-Diferenci ivel]\label{def:weak_D_strategy}
Uma estrat cgia previs -vel $x = (x_t)_{t \geq 0}$  c dita
\textbf{fracamente $D$-diferenci ivel} se, para toda fun SS  o teste suave
$\varphi$ com suporte compacto, o processo $Y_t^\varphi :=
\int_0^\infty \varphi(s)\,x_{t+s}\,ds$  c $D$-diferenci ivel.  A
\textbf{derivada fraca de Nelson} $\mathcal{D}x_t$  c definida pela
rela SS  o:
\begin{equation}\label{eq:weak_D}
  \int_0^\infty \varphi(s)\,\mathcal{D}x_{t+s}\,ds
  := \mathcal{D}\!\left(\int_0^\infty \varphi(s)\,x_{t+s}\,ds\right),
\end{equation}
para toda $\varphi \in C_c^\infty(\mathbb{R}_+)$.
\end{defi}

\begin{obs}[Por que ``fraca''?]
A diferenciabilidade pontual (trajet 3ria a trajet 3ria)  c uma exig ancia
demasiado forte para processos estoc isticos.  A no SS  o fraca (no sentido
das distribui SS ues)  c suficiente para as manipula SS ues anal -ticas que a
GAT requer e  c satisfeita por todas as estrat cgias de interesse pr itico.
Em particular, qualquer estrat cgia adaptada com trajet 3rias c  dl  g de
varia SS  o finita  c fracamente $D$-diferenci ivel.
\end{obs}

% ----------------------------------------------------------------------
\section{Autofinanciamento em Termos de $D$}
% ----------------------------------------------------------------------

Agora reformulamos a condi SS  o de autofinanciamento usando a derivada de
Nelson.  Esta reformula SS  o  c o elo que conecta a teoria cl issica
(Cap -tulo~4) com a formula SS  o geom ctrica da GAT.

\begin{teo}[Autofinanciamento em Termos de $D$ ---
           Proposi SS  o~33 do Artigo]\label{teo:self_financing_D}
Seja $x$ uma estrat cgia fracamente $D$-diferenci ivel e $D_t$ o processo
de deflatores (o valor do instrumento financeiro associado a cada
ativo).  Ent  o:

\begin{enumerate}[(i)]
  \item \textbf{Primeira forma:}
        \begin{equation}\label{eq:sf_D_1}
          \mathcal{D}(x_t \cdot D_t)
          = x_t \cdot \mathcal{D}D_t - \frac{1}{2}\,
            \mathcal{D}\langle x, D\rangle_t.
        \end{equation}

  \item \textbf{Segunda forma (simplificada):}
        \begin{equation}\label{eq:sf_D_2}
          \mathcal{D}x_t \cdot D_t = -\frac{1}{2}\,
          \mathcal{D}\langle x, D\rangle_t.
        \end{equation}
\end{enumerate}
Em ambas as express ues, $\langle x, D\rangle_t$ denota o processo de
covaria SS  o quadr itica (parte cont -nua) entre a estrat cgia $x$ e o
deflator $D$, e $\mathcal{D}\langle x, D\rangle_t$  c sua derivada de
Nelson.
\end{teo}

\begin{proof}
A demonstra SS  o utiliza a regra de Leibniz de Nelson
(Proposi SS  o~\ref{prop:nelson_leibniz}) e a rela SS  o entre It ' e
Stratonovich.

\textbf{Passo 1:} Pela regra de Leibniz de
Nelson~\eqref{eq:nelson_leibniz}, aplicada ao produto escalar $x_t
\cdot D_t = \sum_j x_t^j D_t^j$:
\begin{align}
  \mathcal{D}(x_t \cdot D_t)
  &= (\mathcal{D}x_t) \cdot D_t + x_t \cdot (\mathcal{D}D_t)
     + \frac{1}{2}\,\mathcal{D}\langle x, D\rangle_t. \label{eq:pf1}
\end{align}

\textbf{Passo 2:} A condi SS  o de autofinanciamento na formula SS  o de
Stratonovich (Cap -tulo~4, equa SS  o~\eqref{eq:self_financing_stratonovich} do
Cap -tulo~4)  c:
\[
dV_t^x = x_t \circ dS_t = x_t \circ d(D_t S_t^0).
\]
Em termos do deflator $D_t$, a din  mica de Stratonovich  c governada
pela derivada de Nelson: $\mathcal{D}D_t$  c o drift de Stratonovich de
$D_t$.  O autofinanciamento implica que a varia SS  o do valor da carteira
$(x_t \cdot D_t)S_t^0$ prov cm exclusivamente da varia SS  o de $D_t$
(ponderada por $x_t$) e da capitaliza SS  o por $S_t^0$.

\textbf{Passo 3:} A condi SS  o essencial  c que a estrat cgia $x$ n  o
introduz deriva pr 3pria no valor da carteira --- ou seja, a contribui SS  o
de $\mathcal{D}x_t$ deve ser exatamente cancelada.  Isto leva a:
\[
(\mathcal{D}x_t) \cdot D_t + \frac{1}{2}\,\mathcal{D}\langle x, D\rangle_t = 0,
\]
que  c precisamente a segunda forma~\eqref{eq:sf_D_2}.

Substituindo em~\eqref{eq:pf1}, obtemos:
\[
\mathcal{D}(x_t \cdot D_t)
= -\frac{1}{2}\,\mathcal{D}\langle x, D\rangle_t
  + x_t \cdot \mathcal{D}D_t
  + \frac{1}{2}\,\mathcal{D}\langle x, D\rangle_t
= x_t \cdot \mathcal{D}D_t - \frac{1}{2}\,\mathcal{D}\langle x, D\rangle_t,
\]
que  c a primeira forma~\eqref{eq:sf_D_1}.
\end{proof}

\begin{obs}[Interpreta SS  o Financeira]
A equa SS  o~\eqref{eq:sf_D_2} tem uma interpreta SS  o not ivel: \textbf{a
derivada de Nelson da estrat cgia, contra -da com o deflator,  c
determinada unicamente pela covaria SS  o entre a estrat cgia e o deflator}.
Em outras palavras, a ``velocidade'' com que voc a ajusta sua carteira
($\mathcal{D}x_t$) n  o  c livre --- ela  c restringida pela condi SS  o de
que voc a n  o est i injetando nem retirando capital.  A  onica fonte de
varia SS  o permitida  c a correla SS  o instant  nea entre suas posi SS ues e os
movimentos dos deflatores.
\end{obs}

\begin{obs}[Por que $\frac{1}{2}$?]
O fator $1/2$ na equa SS  o~\eqref{eq:sf_D_2}  c a assinatura da
integral de Stratonovich.  Ele aparece porque a derivada de Nelson
faz a m cdia sim ctrica entre passado e futuro, enquanto a integral de
It ' usa apenas o passado.  Na passagem de It ' para Stratonovich, o
termo de corre SS  o  c exatamente $\frac{1}{2}$ da covaria SS  o quadr itica.
\end{obs}

% ----------------------------------------------------------------------
\section{Conex  o com Malaney--Weinstein}
% ----------------------------------------------------------------------

\subsection{A Conex  o de Malaney--Weinstein}

A ideia de modelar problemas econ 'micos com geometria diferencial n  o
surgiu com Ilinski ou Farinelli.  Malaney (1996)\footnote{Malaney, P.N.
\textit{The Index Number Problem: A Differential Geometry Approach}.
Tese de Doutorado, Harvard University, 1996. \textbf{Resenha:} Primeira
formula SS  o do problema de n omeros- -ndice em linguagem de geometria
diferencial.  Malaney prop 's que a invari  ncia sob mudan SSa de cesta de
bens deveria ser tratada como simetria de gauge.} e Weinstein
(2006)\footnote{Weinstein, E. \textit{Gauge Theory and Inflation}.
Tese de Doutorado, Harvard University, 2006. \textbf{Resenha:} Prop 's
uma teoria de gauge para infla SS  o econ 'mica, tratando  -ndices de pre SSos
como conex ues em fibrados.  Formalizou matematicamente a intui SS  o de
que mudan SSas de unidade de conta s  o transforma SS ues de gauge.}
desenvolveram, em contextos distintos (n omeros- -ndice e infla SS  o), uma
formula SS  o geom ctrica baseada em:

\begin{defi}[Conex  o de Malaney--Weinstein]\label{def:MW_connection}
Em um fibrado sobre o espa SSo de cestas de bens (ou carteiras), a
conex  o de Malaney--Weinstein  c definida pela 1-forma:
\begin{equation}\label{eq:MW_connection}
  \omega^{\text{MW}}(x)(\delta x)
  = \frac{\sum_{j} p_j(x)\,\delta x^j}{\sum_{j} p_j(x)\,x^j},
\end{equation}
onde $p_j(x)$  c o pre SSo do bem $j$ na cesta $x$.  Esta conex  o mede a
varia SS  o relativa do custo da cesta quando a composi SS  o muda.
\end{defi}

\subsection{Equival ancia no Caso Estoc istico}

\begin{prop}[Equival ancia das Conex ues no Contexto da GAT]
\label{prop:MW_equivalence}
A conex  o de Ilinski--Farinelli~\eqref{eq:connection_definition}
reduz-se    conex  o de Malaney--Weinstein quando:
\begin{enumerate}[(i)]
  \item Restringimo-nos    dire SS  o puramente nominal ($\delta t = 0$);
  \item Identificamos os pre SSos $p_j$ com os deflatores $D_t^j$;
  \item O grupo de gauge  c restrito   s transforma SS ues que preservam a
        normaliza SS  o $P_{t,t} = 1$.
\end{enumerate}
Neste limite, ambas as conex ues descrevem a \textbf{taxa de c  mbio
infinitesimal} entre carteiras, e a curvatura de $\omega^{\text{MW}}$
mede a \textbf{inconsist ancia de  -ndices de pre SSos} --- o an ilogo
macroecon 'mico da arbitragem microecon 'mica.
\end{prop}

\begin{proof}[Esbo SSo]
Para $\delta t = 0$, a conex  o de Ilinski--Farinelli  c:
\[
\chi(x,t)(\delta x, 0) = \frac{D_t^{\delta x}}{D_t^x}
= \frac{\sum_j D_t^j\,\delta x^j}{\sum_j D_t^j\,x^j},
\]
que  c exatamente a forma de Malaney--Weinstein com $p_j = D_t^j$.  A
diferen SSa essencial  c que a GAT adiciona a dimens  o temporal (com o
termo $-r_t^x\,\delta t$) e trabalha com o grupo completo $G^*$ (todas
as transforma SS ues de gauge sobre a estrutura a termo), n  o apenas as
transforma SS ues que preservam o ponto $t$.
\end{proof}

\begin{obs}[Generaliza SS  o da GAT]
A conex  o de Ilinski--Farinelli \textbf{cont cm} a conex  o de
Malaney--Weinstein como caso particular, mas a generaliza em tr as
dire SS ues cruciais:
\begin{enumerate}
  \item \textbf{Dimens  o temporal expl -cita:} o termo $-r_t^x\,dt$
        introduz a din  mica de desconto, ausente em Malaney--Weinstein;
  \item \textbf{Estrutura a termo completa:} o grupo $G^*$ age sobre
        toda a estrutura $P_{t,s}$, n  o apenas sobre pre SSos spot;
  \item \textbf{Contexto estoc istico:} a derivada de Nelson e o c ilculo
        de Stratonovich fornecem o arcabou SSo anal -tico para lidar com
        difus ues e saltos, enquanto Malaney--Weinstein  c puramente
        determin -stico.
\end{enumerate}
 %% precisamente esta generaliza SS  o que permite    GAT capturar a
arbitragem como curvatura --- a dimens  o temporal e a estrutura a termo
s  o indispens iveis para que a curvatura tenha significado financeiro.
\end{obs}

% ----------------------------------------------------------------------
\section{Resumo do Cap -tulo}
% ----------------------------------------------------------------------

A conex  o de Ilinski--Farinelli  c o objeto central da GAT.  Neste
cap -tulo, estabelecemos:

\begin{enumerate}
  \item \textbf{Defini SS  o da conex  o:} $\chi(x,t)(\delta x, \delta t)
        = D_t^{\delta x}/D_t^x - r_t^x\,\delta t$, que unifica as
        opera SS ues de c  mbio (ao longo de $x$) e desconto (ao longo de
        $t$) em uma  onica 1-forma com valores na  ilgebra de Lie de
        $G^*$.

  \item \textbf{Forma de conex  o:} em coordenadas,
        $\chi = (1/D_t^x)\sum_j D_t^j\,dx^j - r_t^x\,dt$.

  \item \textbf{Transporte paralelo (Teorema~29):}
        \begin{itemize}
          \item Ao longo de $x$-linhas: multiplica SS  o pela taxa de
                c  mbio $D_{t_0}^{x_0}/D_{t_0}^{x_0 + \Delta x}$.
          \item Ao longo de $t$-linhas: divis  o pelo fator de desconto
                $\exp(\int_{t_0}^{t_0+\Delta t} r_u^{x_0}\,du)$.
        \end{itemize}

  \item \textbf{Derivada de Nelson:} $\mathcal{D}X_t$, a derivada
        sim ctrica que satisfaz a regra da cadeia e  c compat -vel com
        Stratonovich --- o elo entre o c ilculo estoc istico e a
        geometria diferencial.

  \item \textbf{Autofinanciamento em $D$ (Proposi SS  o~33):}
        $\mathcal{D}(x_t \cdot D_t) = x_t \cdot \mathcal{D}D_t
        - \frac{1}{2}\mathcal{D}\langle x, D\rangle_t$, e a forma
        simplificada $\mathcal{D}x_t \cdot D_t
        = -\frac{1}{2}\mathcal{D}\langle x, D\rangle_t$.

  \item \textbf{Conex  o com Malaney--Weinstein:} a GAT generaliza a
        formula SS  o determin -stica original, adicionando a dimens  o
        temporal e a estrutura a termo.
\end{enumerate}

No pr 3ximo cap -tulo, usaremos a conex  o aqui definida para calcular a
\textbf{curvatura} do fibrado do mercado --- e demonstraremos o
resultado central da GAT: a equival ancia entre curvatura zero e aus ancia
de arbitragem.

\begin{center}
\fbox{\begin{minipage}{0.88\textwidth}
\textbf{Mensagem Central do Cap -tulo:} A conex  o de mercado n  o  c um
artefato matem itico ---  c a \textbf{codifica SS  o geom ctrica de todas as
rela SS ues de valor} entre instrumentos financeiros.  Transporte paralelo
ao longo de $x$  c c  mbio; ao longo de $t$  c desconto.  Quando estas
duas opera SS ues s  o consistentes entre si (i.e., quando o transporte n  o
depende do caminho), a curvatura se anula --- e o mercado  c livre de
arbitragem.
\end{minipage}}
\end{center}

% ======================================================================
% REFERENCIAS DO CAPITULO (selecao principal)
% ======================================================================
\begin{thebibliography}{99}

\bibitem{farinelli2021}
Farinelli, S.
\textit{Geometric Arbitrage Theory and Market Dynamics Reloaded}.
arXiv:0910.1671v10, 2021.

\bibitem{ilinski2001}
Ilinski, K.
\textit{Physics of Finance: Gauge Modelling in Non-Equilibrium Pricing}.
Wiley, Chichester, 2001.
ISBN: 978-0471877387.

\bibitem{malaney1996}
Malaney, P.N.
\textit{The Index Number Problem: A Differential Geometry Approach}.
Tese de Doutorado, Harvard University, 1996.

\bibitem{nelson1967}
Nelson, E.
\textit{Dynamical Theories of Brownian Motion}.
Princeton University Press, Princeton, 1967.
ISBN: 978-0691079509.

\bibitem{weinstein2006}
Weinstein, E.
\textit{Gauge Theory and Inflation}.
Tese de Doutorado, Harvard University, 2006.

\bibitem{kobayashi1963}
Kobayashi, S.; Nomizu, K.
\textit{Foundations of Differential Geometry}, Vol.~I.
Wiley, New York, 1963.
ISBN: 978-0471157336.

\bibitem{bleecker1981}
Bleecker, D.
\textit{Gauge Theory and Variational Principles}.
Addison-Wesley, 1981. Reimpresso por Dover, 2005.
ISBN: 978-0486445462.

\end{thebibliography}

% ======================================================================
% FIM DO CAPITULO 7
% ======================================================================


